%%%
%%% Radical "⽶"
%%%
\section*{Radical 119: ``⽶''}\addcontentsline{toc}{section}{Radical 119: ⽶}\addcontentsline{loh}{figure}{\#\#\#\# 119: ⽶}

%%%%%%%%%% 米 %%%%%%%%%%
\subsection*{米}\addcontentsline{loh}{figure}{米}

\begin{Entry}{米}{6}{⽶}[Kangxi 119]
  \begin{Phonetics}{米}{mi3}[][HSK 2,3]
    \definition*{s.}{Sobrenome: Mi}
    \definition{clas.}{m, metro; unidade principal de comprimento do sistema métrico}
    \definition[粒,斤]{s.}{arroz | sementes descascadas; refere"-se a sementes comestíveis descascadas ou sem casca | qualquer coisa que se assemelhe a um grão de arroz}
  \end{Phonetics}
\end{Entry}

\begin{Entry}{米饭}{6,7}{⽶,⾷}
  \begin{Phonetics}{米饭}{mi3fan4}[][HSK 1]
    \definition[碗,盆,锅]{s.}{arroz (cozido)}
  \end{Phonetics}
\end{Entry}

%%%%%%%%%% 类 %%%%%%%%%%
\subsection*{类}\addcontentsline{loh}{figure}{类}

\begin{Entry}{类}{9}{⽶}
  \begin{Phonetics}{类}{lei4}[][HSK 3]
    \definition*{s.}{Sobrenome: Lei}
    \definition{clas.}{tipo; espécie; categoria usada para pessoas ou coisas}
    \definition{s.}{classe; categoria; tipo; variedade; a combinação de muitas coisas semelhantes ou iguais}
    \definition{v.}{assemelhar"-se a; ser semelhante a}
  \synonymref{似}{shi4}
  \synonymref{似}{si4}
  \synonymref{种}{zhong3}
  \end{Phonetics}
\end{Entry}

\begin{Entry}{类似}{9,6}{⽶,⼈}
  \begin{Phonetics}{类似}{lei4si4}[][HSK 3]
    \definition{adj.}{semelhante; análogo}
  \synonymref{仿佛}{fang3fu2}
  \synonymref{好像}{hao3xiang4}
  \synonymref{相似}{xiang1si4}
  \synonymref{相像}{xiang1xiang4}
  \antonymref{相反}{xiang1fan3}
  \end{Phonetics}
\end{Entry}

\begin{Entry}{类别}{9,7}{⽶,⼑}
  \begin{Phonetics}{类别}{lei4bie2}[][HSK 7-9]
    \definition[个,种]{s.}{categoria; classificação; diferentes tipos de coisas, ideias, etc.}
  \end{Phonetics}
\end{Entry}

\begin{Entry}{类型}{9,9}{⽶,⼟}
  \begin{Phonetics}{类型}{lei4xing2}[][HSK 4]
    \definition[种,个]{s.}{tipo; espécie; categoria; tipos formados por coisas com características comuns}
  \synonymref{类别}{lei4bie2}
  \synonymref{品种}{pin3zhong3}
  \synonymref{种类}{zhong3lei4}
  \end{Phonetics}
\end{Entry}

%%%%%%%%%% 粉 %%%%%%%%%%
\subsection*{粉}\addcontentsline{loh}{figure}{粉}

\begin{Entry}{粉}{10}{⽶}
  \begin{Phonetics}{粉}{fen3}[][HSK 7-9]
    \definition{adj.}{branco | rosa}
    \definition{s.}{pó | cosméticos em pó | farinha de trigo | macarrão ou outro alimento feito de feijão, arroz, batata, amido de batata"-doce, etc. | macarrão de arroz}
    \definition{v.}{virar pó | Dialeto: caiar}
  \end{Phonetics}
\end{Entry}

\begin{Entry}{粉丝}{10,5}{⽶,⼀}
  \begin{Phonetics}{粉丝}{fen3si1}[][HSK 7-9]
    \definition[个,群,位,名,些,批]{s.}{aletria de amido de feijão ou batata; aletria chinesa; macarrão de celofane ou macarrão de vidro (transparente) | Empréstimo linguístico: fã; refere"-se a uma pessoa que é obcecada ou adora uma celebridade}
  \end{Phonetics}
\end{Entry}

\begin{Entry}{粉色}{10,6}{⽶,⾊}
  \begin{Phonetics}{粉色}{fen3se4}
    \definition{s.}{cor"-de"-rosa; frequentemente usada para descrever eventos que envolvem romance ou paixão}
  \synonymref{红色}{hong2se4}
  \end{Phonetics}
\end{Entry}

\begin{Entry}{粉碎}{10,13}{⽶,⽯}
  \begin{Phonetics}{粉碎}{fen3sui4}[][HSK 7-9]
    \definition{adj.}{pulverizado; quebrado em pedaços; descreve algo que está muito quebrado, quebrado em partículas muito pequenas}
    \definition{v.}{esmagar; transformar as coisas em partículas muito pequenas | esmagar; quebrar; estilhaçar; fazer com que a outra parte falhe ou seja completamente destruída}
  \end{Phonetics}
\end{Entry}

%%%%%%%%%% 粒 %%%%%%%%%%
\subsection*{粒}\addcontentsline{loh}{figure}{粒}

\begin{Entry}{粒}{11}{⽶}
  \begin{Phonetics}{粒}{li4}[][HSK 7-9]
    \definition{clas.}{utilizado para materiais granulares, grânulos}
    \definition{s.}{partículas pequenas; grão; grânulo; \emph{pellet}}
  \seealsoref{粒儿}{li4r5}
  \end{Phonetics}
\end{Entry}

\begin{Entry}{粒儿}{11,2}{⽶,⼉}
  \begin{Phonetics}{粒儿}{li4r5}
    \definition{s.}{grão}
  \seealsoref{粒}{li4}
  \end{Phonetics}
\end{Entry}

%%%%%%%%%% 粗 %%%%%%%%%%
\subsection*{粗}\addcontentsline{loh}{figure}{粗}

\begin{Entry}{粗}{11}{⽶}
  \begin{Phonetics}{粗}{cu1}[][HSK 4]
    \definition{adj.}{largo (em diâmetro); grosso | grosseiro; rude; áspero | áspero; rouco | descuidado; negligente | rude; sem refinamento; vulgar}
    \definition{adv.}{grosseiramente; vagamente}
  \synonymref{长}{chang2}
  \synonymref{大}{da4}
  \synonymref{宽}{kuan1}
  \synonymref{略}{lve4}
  \synonymref{疏}{shu1}
  \synonymref{俗}{su2}
  \synonymref{微}{wei1}
  \antonymref{短}{duan3}
  \antonymref{精}{jing1}
  \antonymref{细}{xi4}
  \antonymref{雅}{ya3}
  \antonymref{窄}{zhai3}
  \end{Phonetics}
\end{Entry}

\begin{Entry}{粗心}{11,4}{⽶,⼼}
  \begin{Phonetics}{粗心}{cu1xin1}[][HSK 4]
    \definition{adj.}{descuidado; irrefletido; (fazer as coisas) de forma desleixada, sem cuidado}
  \seealsoref{大意}{da4yi5}
  \synonymref{忽略}{hu1lve4}
  \synonymref{忽视}{hu1shi4}
  \synonymref{马虎}{ma3hu5}
  \synonymref{疏忽}{shu1hu5}
  \antonymref{谨慎}{jin3shen4}
  \antonymref{精心}{jing1xin1}
  \antonymref{留神}{liu2//shen2}
  \antonymref{认真}{ren4zhen1}
  \antonymref{细心}{xi4xin1}
  \antonymref{细致}{xi4zhi4}
  \antonymref{小心}{xiao3xin5}
  \antonymref{严谨}{yan2jin3}
  \antonymref{仔细}{zi3xi4}
  \end{Phonetics}
\end{Entry}

\begin{Entry}{粗心大意}{11,4,3,13}{⽶,⼼,⼤,⼼}
  \begin{Phonetics}{粗心大意}{cu1xin1-da4yi4}[][HSK 7-9]
    \definition{expr.}{ser negligente; descuidado; inadvertido; desmiolado; descuidado e negligente; negligente; remisso; refere"-se a fazer as coisas de forma descuidada}
  \end{Phonetics}
\end{Entry}

\begin{Entry}{粗心地做}{11,4,6,11}{⽶,⼼,⼟,⼈}
  \begin{Phonetics}{粗心地做}{cu1xin1di4 zuo4}
    \definition{adj.}{feito descuidadamente; embaralhado}
  \end{Phonetics}
\end{Entry}

\begin{Entry}{粗略}{11,11}{⽶,⽥}
  \begin{Phonetics}{粗略}{cu1lve4}[][HSK 7-9]
    \definition{adj.}{grosseiro; rudimentar; superficial}
  \end{Phonetics}
\end{Entry}

\begin{Entry}{粗鲁}{11,12}{⽶,⿂}
  \begin{Phonetics}{粗鲁}{cu1lu3}[][HSK 7-9]
    \definition{adj.}{rude; grosseiro; incivilizado}
  \end{Phonetics}
\end{Entry}

\begin{Entry}{粗暴}{11,15}{⽶,⽇}
  \begin{Phonetics}{粗暴}{cu1bao4}[][HSK 7-9]
    \definition{adj.}{rude; áspero; bruto; brutal; violento}
  \end{Phonetics}
\end{Entry}

\begin{Entry}{粗糙}{11,16}{⽶,⽶}
  \begin{Phonetics}{粗糙}{cu1cao1}[][HSK 7-9]
    \definition{adj.}{áspero; grosseiro; não é liso; não é redondo; não é fino | desleixado; descuidado; não meticuloso}
  \end{Phonetics}
\end{Entry}

%%%%%%%%%% 粘 %%%%%%%%%%
\subsection*{粘}\addcontentsline{loh}{figure}{粘}

\begin{Entry}{粘}{11}{⽶}
  \begin{Phonetics}{粘}{nian2}
    \variantof{黏}
  \end{Phonetics}
  \begin{Phonetics}{粘}{zhan1}[][HSK 7-9]
    \definition{v.}{colar; grudar; afixar; unir; aderir substâncias pegajosas a objetos ou umas às outras; usar adesivo para unir objetos}
  \end{Phonetics}
\end{Entry}

%%%%%%%%%% 粥 %%%%%%%%%%
\subsection*{粥}\addcontentsline{loh}{figure}{粥}

\begin{Entry}{粥}{12}{⽶}
  \begin{Phonetics}{粥}{yu4}
    \definition{v.}{dar a luz; ter filhos}
  \end{Phonetics}
  \begin{Phonetics}{粥}{zhou1}[][HSK 6]
    \definition[碗,锅,口]{s.}{mingau; mingau de aveia; alimentos semilíquidos feitos de grãos ou grãos misturados com outras coisas}
  \end{Phonetics}
\end{Entry}

%%%%%%%%%% 粪 %%%%%%%%%%
\subsection*{粪}\addcontentsline{loh}{figure}{粪}

\begin{Entry}{粪}{12}{⽶}
  \begin{Phonetics}{粪}{fen4}[][HSK 7-9]
    \definition[缸,桶]{s.}{excremento; fezes; esterco}
    \definition{v.}{Literário: aplicar esterco; fertilizar | Literário: limpar; remover; eliminar; acabar com}
  \end{Phonetics}
\end{Entry}

\begin{Entry}{粪便}{12,9}{⽶,⼈}
  \begin{Phonetics}{粪便}{fen4bian4}[][HSK 7-9]
    \definition{s.}{excremento; fezes; esterco; excrementos e urina}
  \end{Phonetics}
\end{Entry}

%%%%%%%%%% 缓 %%%%%%%%%%
\subsection*{缓}\addcontentsline{loh}{figure}{缓}

\begin{Entry}{缓}{12}{⽶}
  \begin{Phonetics}{缓}{huan3}[][HSK 7-9]
    \definition{adj.}{lento; sem pressa | sem tensão; relaxado}
    \definition{v.}{atrasar; adiar; protelar | recuperar; reviver; voltar a si}
  \end{Phonetics}
\end{Entry}

\begin{Entry}{缓和}{12,8}{⽶,⼝}
  \begin{Phonetics}{缓和}{huan3he2}[][HSK 7-9]
    \definition{adj.}{relaxado; moderado; suave; pacífico e relaxante; não tenso ou intenso}
    \definition{v.}{relaxar; aliviar; atenuar; facilitar}
  \end{Phonetics}
\end{Entry}

\begin{Entry}{缓缓}{12,12}{⽶,⽶}
  \begin{Phonetics}{缓缓}{huan3huan3}[][HSK 7-9]
    \definition{adv.}{lentamente; vagarosamente; gradualmente}
  \end{Phonetics}
\end{Entry}

\begin{Entry}{缓解}{12,13}{⽶,⾓}
  \begin{Phonetics}{缓解}{huan3jie3}[][HSK 4]
    \definition{v.}{facilitar; aliviar; atenuar; amenizar; reduzir}
  \synonymref{缓和}{huan3he2}
  \synonymref{减轻}{jian3qing1}
  \antonymref{恶化}{e4hua4}
  \antonymref{激化}{ji1hua4}
  \antonymref{加剧}{jia1ju4}
  \antonymref{加深}{jia1shen1}
  \antonymref{加重}{jia1zhong4}
  \antonymref{升级}{sheng1//ji2}
  \end{Phonetics}
\end{Entry}

\begin{Entry}{缓慢}{12,14}{⽶,⼼}
  \begin{Phonetics}{缓慢}{huan3man4}[][HSK 7-9]
    \definition{adj.}{lento; vagaroso}
    \definition{adv.}{lentamente; vagarosamente}
  \end{Phonetics}
\end{Entry}

%%%%%%%%%% 粮 %%%%%%%%%%
\subsection*{粮}\addcontentsline{loh}{figure}{粮}

\begin{Entry}{粮}{13}{⽶}
  \begin{Phonetics}{粮}{liang2}
    \definition[斤,粒]{s.}{grãos; alimentos; provisões | imposto sobre grãos | nutrição | imposto agrícola; grãos como imposto agrícola}
  \end{Phonetics}
\end{Entry}

\begin{Entry}{粮食}{13,9}{⽶,⾷}
  \begin{Phonetics}{粮食}{liang2shi5}[][HSK 4]
    \definition[种,吨,袋,颗,粒]{s.}{alimentos; grãos; termo geral para os vários tipos de arroz, feijão, etc. que podem ser consumidos}
  \end{Phonetics}
\end{Entry}

%%%%%%%%%% 粽 %%%%%%%%%%
\subsection*{粽}\addcontentsline{loh}{figure}{粽}

\begin{Entry}{粽}{14}{⽶}
  \begin{Phonetics}{粽}{zong4}
    \definition[个]{s.}{bolinho em forma de pirâmide feito de arroz glutinoso envolto em folhas de bambu ou junco (geralmente consumido durante o Festival do Barco"-Dragão, 端午节) | bolinhos de arroz envoltos em folhas}
  \seealsoref{端午节}{duan1wu3jie2}
  \seealsoref{粽子}{zong4zi5}
  \end{Phonetics}
\end{Entry}

\begin{Entry}{粽子}{14,3}{⽶,⼦}
  \begin{Phonetics}{粽子}{zong4zi5}[][HSK 7-9]
    \definition[个,只,颗,锅]{s.}{arroz glutinoso com recheio à escolha, envolto em folhas e cozido; o zongzi é um tipo de alimento feito com arroz glutinoso enrolado em folhas de bambu ou junco, moldado em um cone triangular ou outro formato, e cozido antes de ser consumido; comer zongzi é um costume tradicional chinês durante o Festival do Barco"-Dragão, 端午节}
  \seealsoref{端午节}{duan1wu3jie2}
  \end{Phonetics}
\end{Entry}

%%%%%%%%%% 精 %%%%%%%%%%
\subsection*{精}\addcontentsline{loh}{figure}{精}

\begin{Entry}{精}{14}{⽶}
  \begin{Phonetics}{精}{jing1}[][HSK 6]
    \definition{adj.}{refinado; escolhido; purificado ou selecionado | perfeito; excelente; melhor | fino; preciso; meticuloso | inteligente; astuto; esperto | habilidoso; versado; proficiente}
    \definition{adv.}{muito; extremamente; antes de certos adjetivos, significa 十分 ou 非常}
    \definition{s.}{extrato; essência; essência refinada ou selecionada; extraída | energia; espírito | semente; esperma; sêmen | \emph{goblin}; espírito; elfo; demônio}
  \seealsoref{非常}{fei1chang2}
  \seealsoref{十分}{shi2fen1}
  \synonymref{好}{hao3}
  \synonymref{美}{mei3}
  \synonymref{细}{xi4}
  \synonymref{妖}{yao1}
  \antonymref{粗}{cu1}
  \antonymref{傻}{sha3}
  \end{Phonetics}
\end{Entry}

\begin{Entry}{精力}{14,2}{⽶,⼒}
  \begin{Phonetics}{精力}{jing1li4}[][HSK 4]
    \definition[些]{s.}{energia; vigor; força mental e física}
  \seealsoref{精神}{jing1shen5}
  \synonymref{活力}{huo2li4}
  \synonymref{生气}{sheng1//qi4}
  \synonymref{生机}{sheng1ji1}
  \synonymref{元气}{yuan2qi4}
  \antonymref{疲惫}{pi2bei4}
  \antonymref{衰弱}{shuai1ruo4}
  \antonymref{无力}{wu2li4}
  \end{Phonetics}
\end{Entry}

\begin{Entry}{精子}{14,3}{⽶,⼦}
  \begin{Phonetics}{精子}{jing1zi3}
    \definition{s.}{espermatozoide; célula germinativa}
  \end{Phonetics}
\end{Entry}

\begin{Entry}{精心}{14,4}{⽶,⼼}
  \begin{Phonetics}{精心}{jing1xin1}[][HSK 7-9]
    \definition{adv.}{meticulosamente; cuidadosamente; elaboradamente; preste muita atenção; concentre"-se totalmente}[他们精心设计了这个项目。===Eles planejaram este projeto meticulosamente.]
  \end{Phonetics}
\end{Entry}

\begin{Entry}{精打细算}{14,5,8,14}{⽶,⼿,⽷,⽵}
  \begin{Phonetics}{精打细算}{jing1da3-xi4suan4}[][HSK 7-9]
    \definition{expr.}{seja muito cuidadoso nos cálculos; contagem precisa; seja preciso nos cálculos; faça um orçamento rigoroso; cálculo cuidadoso e detalhado (meticuloso); cálculo cuidadoso e orçamento rigoroso; conte cada centavo e faça cada centavo valer a pena; corte os gastos com precisão; planejar com meticulosidade e cuidado, isso significa calcular com precisão o uso de mão de obra, recursos materiais e recursos financeiros para evitar desperdícios}
  \end{Phonetics}
\end{Entry}

\begin{Entry}{精华}{14,6}{⽶,⼗}
  \begin{Phonetics}{精华}{jing1hua2}[][HSK 7-9]
    \definition{s.}{elite; creme; escolha; essência; quintessência; a melhor e mais refinada parte de tudo | glória; esplendor; brilho; luz (do Sol e da Lua)}
  \end{Phonetics}
\end{Entry}

\begin{Entry}{精妙}{14,7}{⽶,⼥}
  \begin{Phonetics}{精妙}{jing1miao4}[][HSK 7-9]
    \definition{adj.}{requintado | fino e delicado (geralmente de obras de arte)}
  \end{Phonetics}
\end{Entry}

\begin{Entry}{精灵}{14,7}{⽶,⽕}
  \begin{Phonetics}{精灵}{jing1ling2}
    \definition{adj.}{esperto; inteligente; perspicaz}
    \definition{s.}{espírito; demônio; fada; elfo; duende; gênio}
  \synonymref{机灵}{ji1ling5}
  \end{Phonetics}
\end{Entry}

\begin{Entry}{精明}{14,8}{⽶,⽇}
  \begin{Phonetics}{精明}{jing1ming2}[][HSK 7-9]
    \definition{adj.}{astuto; sagaz; perspicaz; inteligente e brilhante}
  \end{Phonetics}
\end{Entry}

\begin{Entry}{精练}{14,8}{⽶,⽷}
  \begin{Phonetics}{精练}{jing1lian4}[][HSK 7-9]
    \definition{adj.}{conciso; sucinto; lacônico | refinado; palavras e frases redundantes eliminadas}
    \definition{v.}{praticar intensivamente}
  \synonymref{干脆}{gan1cui4}
  \synonymref{简单}{jian3dan1}
  \synonymref{简洁}{jian3jie2}
  \synonymref{精炼}{jing1lian4}
  \antonymref{冗长}{rong3chang2}
  \end{Phonetics}
\end{Entry}

\begin{Entry}{精细}{14,8}{⽶,⽷}
  \begin{Phonetics}{精细}{jing1xi4}[][HSK 7-9]
    \definition{adj.}{fino; cuidadoso; meticuloso; muito delicado | astuto; perspicaz e cuidadoso; muito meticuloso}
  \end{Phonetics}
\end{Entry}

\begin{Entry}{精英}{14,8}{⽶,⾋}
  \begin{Phonetics}{精英}{jing1ying1}[][HSK 7-9]
    \definition{s.}{creme; essência; quintessência | escolhido; elite; pessoa de habilidade excepcional}
  \end{Phonetics}
\end{Entry}

\begin{Entry}{精品}{14,9}{⽶,⼝}
  \begin{Phonetics}{精品}{jing1pin3}[][HSK 6]
    \definition[个]{s.}{belas obras (de arte); objetos de arte | produtos de qualidade; artigos de excelente qualidade; produto \emph{premium}}
  \synonymref{珍藏}{zhen1cang2}
  \antonymref{废品}{fei4pin3}
  \antonymref{垃圾}{la1ji1}
  \end{Phonetics}
\end{Entry}

\begin{Entry}{精炼}{14,9}{⽶,⽕}
  \begin{Phonetics}{精炼}{jing1lian4}
    \definition{adj.}{conciso; sucinto; lacônico}
    \definition{s.}{refino; remoção de impurezas}
    \definition{v.}{Metalurgia: refinar; purificar; fundir}
  \end{Phonetics}
\end{Entry}

\begin{Entry}{精神}{14,9}{⽶,⽰}
  \begin{Phonetics}{精神}{jing1shen2}[][HSK 3]
    \definition[种,个,类,股]{s.}{espírito; mente; estado mental; refere"-se à consciência, às atividades mentais e ao estado psicológico geral de uma pessoa | substância; espírito; essência; propósito; significado principal}
  \seealsoref{精力}{jing1li4}
  \synonymref{思想}{si1xiang3}
  \synonymref{心灵}{xin1ling2}
  \synonymref{意识}{yi4shi5}
  \synonymref{元气}{yuan2qi4}
  \synonymref{宗旨}{zong1zhi3}
  \antonymref{身体}{shen1ti3}
  \antonymref{物质}{wu4zhi4}
  \end{Phonetics}
  \begin{Phonetics}{精神}{jing1shen5}[][HSK 3]
    \definition{adj.}{animado; espirituoso; vigoroso; descreve uma pessoa como cheia de energia | muito bonito; boa aparência, bom físico}
    \definition[种,个,类,股]{s.}{impulso; vigor; vitalidade}
  \seealsoref{精力}{jing1li4}
  \synonymref{活力}{huo2li4}
  \synonymref{生气}{sheng1//qi4}
  \synonymref{英俊}{ying1jun4}
  \synonymref{元气}{yuan2qi4}
  \antonymref{疲惫}{pi2bei4}
  \antonymref{疲倦}{pi2juan4}
  \antonymref{疲劳}{pi2lao2}
  \antonymref{物质}{wu4zhi4}
  \end{Phonetics}
\end{Entry}

\begin{Entry}{精神病}{14,9,10}{⽶,⽰,⽧}
  \begin{Phonetics}{精神病}{jing1shen2bing4}[][HSK 7-9]
    \definition{s.}{doença mental; transtorno mental; psicose}[这是幻想型精神病的体现。===Isso é uma manifestação de psicose delirante.]
  \end{Phonetics}
\end{Entry}

\begin{Entry}{精美}{14,9}{⽶,⽺}
  \begin{Phonetics}{精美}{jing1mei3}[][HSK 6]
    \definition{adj.}{elegante; requintado}
  \seealsoref{精致}{jing1zhi4}
  \synonymref{精妙}{jing1miao4}
  \synonymref{精细}{jing1xi4}
  \synonymref{优美}{you1mei3}
  \antonymref{粗糙}{cu1cao1}
  \antonymref{简陋}{jian3lou4}
  \antonymref{劣质}{lie4zhi4}
  \end{Phonetics}
\end{Entry}

\begin{Entry}{精疲力竭}{14,10,2,14}{⽶,⽧,⼒,⽴}
  \begin{Phonetics}{精疲力竭}{jing1pi2-li4jie2}[][HSK 7-9]
    \definition{expr.}{esgotado; exausto; desgastado; descrevendo fadiga extrema e completa falta de energia}
  \end{Phonetics}
\end{Entry}

\begin{Entry}{精益求精}{14,10,7,14}{⽶,⽫,⽔,⽶}
  \begin{Phonetics}{精益求精}{jing1yi4qiu2jing1}[][HSK 7-9]
    \definition{expr.}{``Busque a excelência.''; esforçar"-se pela perfeição; buscar a melhoria constante; perseguir a excelência; almejar a perfeição; já está muito bom, mas você ainda quer que fique ainda melhor; melhorar algo constantemente; continuar melhorando}
  \end{Phonetics}
\end{Entry}

\begin{Entry}{精致}{14,10}{⽶,⾄}
  \begin{Phonetics}{精致}{jing1zhi4}[][HSK 7-9]
    \definition{adj.}{fino; requintado; delicado}[我们欣赏她精致的手工艺品。===Admiramos seu trabalho artesanal requintado.]
  \end{Phonetics}
\end{Entry}

\begin{Entry}{精通}{14,10}{⽶,⾡}
  \begin{Phonetics}{精通}{jing1tong1}[][HSK 7-9]
    \definition{v.}{dominar; ser proficiente em; ter um bom domínio de; ter um profundo entendimento e conhecimento abrangente de uma área específica de estudo, tecnologia ou negócios}
  \end{Phonetics}
\end{Entry}

\begin{Entry}{精密}{14,11}{⽶,⼧}
  \begin{Phonetics}{精密}{jing1mi4}
    \definition{adj.}{preciso; preciso e meticuloso}
  \end{Phonetics}
\end{Entry}

\begin{Entry}{精彩}{14,11}{⽶,⼺}
  \begin{Phonetics}{精彩}{jing1cai3}[][HSK 3]
    \definition{adj.}{brilhante; esplêndido; maravilhoso}
  \seealsoref{漂亮}{piao4liang5}
  \synonymref{出色}{chu1se4}
  \synonymref{精华}{jing1hua2}
  \synonymref{精美}{jing1mei3}
  \synonymref{美妙}{mei3miao4}
  \synonymref{卓越}{zhuo2yue4}
  \antonymref{平淡}{ping2dan4}
  \antonymref{平庸}{ping2yong1}
  \antonymref{无聊}{wu2liao2}
  \antonymref{糟糕}{zao1gao1}
  \antonymref{拙劣}{zhuo1lie4}
  \end{Phonetics}
\end{Entry}

\begin{Entry}{精确}{14,12}{⽶,⽯}
  \begin{Phonetics}{精确}{jing1que4}[][HSK 7-9]
    \definition{adj.}{exato; preciso; acurado; muito preciso e correto}
  \end{Phonetics}
\end{Entry}

\begin{Entry}{精简}{14,13}{⽶,⽵}
  \begin{Phonetics}{精简}{jing1jian3}[][HSK 7-9]
    \definition{v.}{reduzir; simplificar; cortar; simplificar; eliminar o desnecessário e conservar o necessário}
  \end{Phonetics}
\end{Entry}

\begin{Entry}{精髓}{14,21}{⽶,⾻}
  \begin{Phonetics}{精髓}{jing1sui3}[][HSK 7-9]
    \definition{s.}{medula; medula óssea; quintessência; essência metafórica das coisas}
  \end{Phonetics}
\end{Entry}

%%%%%%%%%% 糆 %%%%%%%%%%
\subsection*{糆}\addcontentsline{loh}{figure}{糆}

\begin{Entry}{糆}{15}{⽶}
  \begin{Phonetics}{糆}{mian4}
    \variantof{面}
  \end{Phonetics}
\end{Entry}

%%%%%%%%%% 糊 %%%%%%%%%%
\subsection*{糊}\addcontentsline{loh}{figure}{糊}

\begin{Entry}{糊}{15}{⽶}
  \begin{Phonetics}{糊}{hu1}
    \definition{v.}{colar; untar; usar uma pasta mais espessa para revestir costuras, furos ou superfícies planas}
  \end{Phonetics}
  \begin{Phonetics}{糊}{hu2}[][HSK 7-9]
    \definition{adj.}{queimado}
    \definition{s.}{mingau; pasta; papa}
    \definition{v.}{colar com pasta; colar | (comida) ser queimado}
  \end{Phonetics}
  \begin{Phonetics}{糊}{hu4}
    \definition{s.}{pasta; comida que parece mingau}
  \end{Phonetics}
\end{Entry}

\begin{Entry}{糊里糊涂}{15,7,15,10}{⽶,⾥,⽶,⽔}
  \begin{Phonetics}{糊里糊涂}{hu2li5hu2tu5}
    \definition{adj.}{desnorteado | perturbado}
  \end{Phonetics}
\end{Entry}

\begin{Entry}{糊涂}{15,10}{⽶,⽔}
  \begin{Phonetics}{糊涂}{hu2tu5}[][HSK 7-9]
    \definition{adj.}{confuso; perplexo; desnorteado; com compreensão pouco clara ou confusa das coisas | confuso; com conteúdo confuso}
  \end{Phonetics}
\end{Entry}

%%%%%%%%%% 糕 %%%%%%%%%%
\subsection*{糕}\addcontentsline{loh}{figure}{糕}

\begin{Entry}{糕}{16}{⽶}
  \begin{Phonetics}{糕}{gao1}
    \definition{s.}{bolo; alimentos feitos de farinha de arroz, farinha de trigo, etc.}
  \end{Phonetics}
\end{Entry}

\begin{Entry}{糕点}{16,9}{⽶,⽕}
  \begin{Phonetics}{糕点}{gao1dian3}
    \definition[名]{s.}{bolo; doce; pastel}
  \synonymref{点心}{dian3xin5}
  \synonymref{面包}{mian4bao1}
  \end{Phonetics}
\end{Entry}

\begin{Entry}{糕点师}{16,9,6}{⽶,⽕,⼱}
  \begin{Phonetics}{糕点师}{gao1dian3 shi1}
    \definition{s.}{confeiteiro}
  \end{Phonetics}
\end{Entry}

\begin{Entry}{糕点店}{16,9,8}{⽶,⽕,⼴}
  \begin{Phonetics}{糕点店}{gao1dian3dian4}
    \definition{s.}{confeitaria}
  \end{Phonetics}
\end{Entry}

%%%%%%%%%% 糖 %%%%%%%%%%
\subsection*{糖}\addcontentsline{loh}{figure}{糖}

\begin{Entry}{糖}{16}{⽶}
  \begin{Phonetics}{糖}{tang2}[][HSK 3]
    \definition[包,斤,勺,袋,块]{s.}{açúcar; um tipo de açúcar; um tipo de composto orgânico, que pode ser dividido em três tipos: monossacarídeos, dissacarídeos e polissacarídeos; é a principal substância que produz energia térmica no corpo humano, como glicose, sacarose, lactose, amido, etc. | açúcar; açúcar comestível; termo geral para açúcar | doces; balas | carboidrato; algo doce e calórico}
  \end{Phonetics}
\end{Entry}

\begin{Entry}{糖尿病}{16,7,10}{⽶,⼫,⽧}
  \begin{Phonetics}{糖尿病}{tang2niao4bing4}[][HSK 7-9]
    \definition{s.}{diabetes; diabetes mellitus; doença crônica causada pela secreção insuficiente de insulina, levando a distúrbios no metabolismo da glicose e níveis elevados de açúcar no sangue}[糖尿病会带来严重的问题。===O diabetes pode causar problemas graves.]
  \end{Phonetics}
\end{Entry}

\begin{Entry}{糖果}{16,8}{⽶,⽊}
  \begin{Phonetics}{糖果}{tang2guo3}[][HSK 7-9]
    \definition[个,颗,包,袋]{s.}{doce; alimentos açucarados, que geralmente contêm suco de frutas, especiarias, leite ou café, etc.}
  \end{Phonetics}
\end{Entry}

\begin{Entry}{糖醋鱼}{16,15,8}{⽶,⾣,⿂}
  \begin{Phonetics}{糖醋鱼}{tang2cu4yu2}
    \definition{s.}{Prato: peixe guisado em molho agridoce; prato tradicional chinês, o peixe é frito até ficar dourado e crocante, sendo servido com um molho agridoce, resultando num delicioso sabor agridoce.}
  \end{Phonetics}
\end{Entry}

%%%%%%%%%% 糟 %%%%%%%%%%
\subsection*{糟}\addcontentsline{loh}{figure}{糟}

\begin{Entry}{糟}{17}{⽶}
  \begin{Phonetics}{糟}{zao1}[][HSK 5]
    \definition{adj.}{pobre; apodrecido; deteriorado | estragado; em uma bagunça; em um estado miserável (terrível) | (situação ou circunstância) ruim; desfavorável}
    \definition{s.}{resíduos de destilação de bebidas alcoólicas; resíduos do processo de fermentação do vinho}
    \definition{v.}{marinar alimentos em vinho ou mosto | desperdiçar; estragar; destruir}
  \synonymref{差}{cha4}
  \synonymref{费}{fei4}
  \synonymref{腐}{fu3}
  \synonymref{坏}{huai4}
  \synonymref{烂}{lan4}
  \synonymref{损}{sun3}
  \antonymref{好}{hao3}
  \end{Phonetics}
\end{Entry}

\begin{Entry}{糟糕}{17,16}{⽶,⽶}
  \begin{Phonetics}{糟糕}{zao1gao1}[][HSK 5]
    \definition{adj.}{(corpo, situação, etc.) muito ruim, péssimo}
    \definition{interj.}{``Que terrível!''; ``Que má sorte!''; ``Muito ruim!''}
  \seealsoref{倒霉}{dao3//mei2}
  \synonymref{恶劣}{e4lie4}
  \antonymref{出色}{chu1se4}
  \antonymref{精彩}{jing1cai3}
  \antonymref{美好}{mei3hao3}
  \antonymref{漂亮}{piao4liang5}
  \end{Phonetics}
\end{Entry}

%%%%% EOF %%%%%

